\section{Zariski decomposition}


% \subsection{On surfaces}



% \subsection[Nakayama-Zariski decomposition]{Nakayama-Zariski $\sigma$-decomposition}


\subsection{Classical Zariski Decomposition}

  In this subsection, let $X$ be a smooth projective surface over $\kkk$.

  The following result is due to Zariski in the effective case and to Fujita in the general case.

  \begin{theorem}
    Let $X$ be a smooth projective surface and $D$ be a pseudo-effective $\bbR$-divisor on $X$. 
    Then there exists a unique decomposition
    \[ D = P + N \]
    with $P$ being nef $\bbR$-divisor and $N = \sum_i a_i E_i$ being effective $\bbR$-divisor 
    such that 
    \[ P \cdot E_i = 0, \quad \text{and} \quad (E_i \cdot E_j)_{i,j} \text{ is negative definite}. \]
    Moreover, if $D$ is a $\bbQ$-divisor, then $P$ and $N$ are also $\bbQ$-divisors.
    We call this decomposition the \emph{Zariski decomposition} of $D$.
  \end{theorem}
  \begin{proof}
    \Yang{}
  \end{proof}

  \begin{example}
    Let $X = \Bl_p \bbP^2$ be the blow-up of $\bbP^2$ at a point $p$, and let $E$ be the exceptional divisor.
      \Yang{}
  \end{example}

  \begin{example}
    Let $X = \bbP_C(\calO \oplus \calL)$ be 
    \Yang{}
  \end{example}

  \begin{proposition}
    Let $D$ be a pseudo-effective $\bbR$-divisor on a smooth projective surface $X$ with Zariski decomposition $D = P + N$. Then we have
    \[ H^0(X, \lfloor mD \rfloor) = H^0(X, \lfloor mP \rfloor) \]
    for any integer $m \geq 0$. In particular, we have $\kappa(X, D) = \kappa(X, P)$.
   \Yang{}
  \end{proposition}

  \begin{corollary}
    Let $D$ be a pseudo-effective $\bbR$-divisor with Zariski decomposition $D = P + N$.
    Then $\vol(D) = P^2$.
   \Yang{}
  \end{corollary}


  \begin{proposition}
    In general, there does not exist a Zariski decomposition for pseudo-effective divisors on higher dimensional varieties.
   \Yang{}
  \end{proposition}

  \begin{example}
    \Yang{}
  \end{example}


\subsection{Nakayama Zariski Decomposition}

  Let $X$ be a normal projective variety and $D$ be a pseudo-effective $\bbR$-divisor on $X$. For any prime divisor $E$ on $X$, we define
  \[ \sigma_E(D) := \inf\ \{ \Mult_E D' \mid D' \sim_\bbR D, D' \geq 0 \}. \]
  Then we define
  \[ N_\sigma(D) := \sum_E \sigma_E(D) E, \quad P_\sigma(D) := D - N_\sigma(D). \]


  \begin{theorem}
    There is a unique decomposition
    \[ D = P_\sigma(D) + N_\sigma(D) \]
    such that $P_\sigma(D)$ is pseudo-effective and $\sigma_E(P_\sigma(D)) = 0$ for any prime divisor $E$ on $X$. 
    We call this decomposition the \emph{Nakayama Zariski decomposition} of $D$.

    \Yang{}
  \end{theorem}


